% Bibliografía manual, ordenada alfabéticamente.
\chapter{Bibliografía}
\label{chap:bibliografia}

\begin{list}{}{%
  \setlength{\leftmargin}{2em}%
  \setlength{\itemindent}{-2em}%
  \setlength{\itemsep}{0.55em}%
  \setlength{\parsep}{0pt}%
}

\item Amihud, Y. (2002). Illiquidity and Stock Returns: Cross-Section and Time-Series Effects.
  \textit{Journal of Financial Markets}, 5(1), 31--56.
  \url{https://doi.org/10.1016/S1386-4181(01)00024-6}

\item Ang, A., Hodrick, R. J., Xing, Y., y Zhang, X. (2006). The Cross-Section of Volatility and
  Expected Returns. \textit{The Journal of Finance}, 61(1), 259--299.
  \url{https://doi.org/10.1111/j.1540-6261.2006.00836.x}

\item Arnott, R. D., Harvey, C. R., y Markowitz, H. (2019). A Backtesting Protocol in the Era of
  Machine Learning. \textit{The Journal of Financial Data Science}, 1(1), 64--74.
  \url{https://doi.org/10.3905/jfds.2019.1.064}

\item Asness, C. S., Frazzini, A., y Pedersen, L. H. (2019). Quality Minus Junk.
  \textit{Review of Accounting Studies}, 24(1), 34--112.
  \url{https://doi.org/10.1007/s11142-018-9470-2}

\item Bailey, D. H., y López de Prado, M. (2014). The Deflated Sharpe Ratio: Correcting for
  Selection Bias, Backtest Overfitting, and Non-Normality.
  \textit{The Journal of Portfolio Management}, 40(5), 94--107.
  \url{https://doi.org/10.3905/jpm.2014.40.5.094}

\item Bailey, D. H., Borwein, J. M., López de Prado, M., y Zhu, Q. J. (2017). The Probability of
  Backtest Overfitting. \textit{Journal of Computational Finance}, 20(4), 39--69.
  \url{https://doi.org/10.21314/JCF.2017.322}

\item Banz, R. W. (1981). The Relationship between Return and Market Value of Common Stocks.
  \textit{Journal of Financial Economics}, 9(1), 3--18.
  \url{https://doi.org/10.1016/0304-405X(81)90018-0}

\item Brown, S. J., Goetzmann, W., Ibbotson, R. G., y Ross, S. A. (1992). Survivorship Bias in
  Performance Studies. \textit{The Review of Financial Studies}, 5(4), 553--580.
  \url{https://doi.org/10.1093/rfs/5.4.553}

\item Carhart, M. M. (1997). On Persistence in Mutual Fund Performance.
  \textit{The Journal of Finance}, 52(1), 57--82.
  \url{https://doi.org/10.1111/j.1540-6261.1997.tb03808.x}

\item Clarke, R., de Silva, H., y Thorley, S. (2002). Portfolio Constraints and the Fundamental Law
  of Active Management. \textit{Financial Analysts Journal}, 58(5), 48--66.

\item Fama, E. F., y French, K. R. (1992). The Cross-Section of Expected Stock Returns.
  \textit{The Journal of Finance}, 47(2), 427--465.
  \url{https://doi.org/10.1111/j.1540-6261.1992.tb04398.x}

\item Fama, E. F., y French, K. R. (1993). Common Risk Factors in the Returns on Stocks and Bonds.
  \textit{Journal of Financial Economics}, 33(1), 3--56.
  \url{https://doi.org/10.1016/0304-405X(93)90023-5}

\item Fama, E. F., y French, K. R. (2015). A Five-Factor Asset Pricing Model.
  \textit{Journal of Financial Economics}, 116(1), 1--22.
  \url{https://doi.org/10.1016/j.jfineco.2014.10.010}

\item Feng, G., Giglio, S., y Xiu, D. (2020). Taming the Factor Zoo: A Test of New Factors.
  \textit{The Journal of Finance}, 75(3), 1327--1370.
  \url{https://doi.org/10.1111/jofi.12883}

\item Frazzini, A., y Pedersen, L. H. (2014). Betting Against Beta.
  \textit{Journal of Financial Economics}, 111(1), 1--25.
  \url{https://doi.org/10.1016/j.jfineco.2013.10.005}

\item Freyberger, J., Neuhierl, A., y Weber, M. (2020). Dissecting Characteristics Nonparametrically.
  \textit{The Review of Financial Studies}, 33(5), 2326--2377.
  \url{https://doi.org/10.1093/rfs/hhz123}

\item Grinold, R. C. (1989). The Fundamental Law of Active Management.
  \textit{The Journal of Portfolio Management}, 15(3), 30--37.
  \url{https://doi.org/10.3905/jpm.1989.409211}

\item Gu, S., Kelly, B., y Xiu, D. (2020). Empirical Asset Pricing via Machine Learning.
  \textit{The Review of Financial Studies}, 33(5), 2223--2273.
  \url{https://doi.org/10.1093/rfs/hhaa009}

\item Hansen, P. R. (2005). A Test for Superior Predictive Ability.
  \textit{Journal of Business \& Economic Statistics}, 23(4), 365--380.
  \url{https://doi.org/10.1198/073500105000000063}

\item Harvey, C. R., Liu, Y., y Zhu, H. (2016). \ldots{} and the Cross-Section of Expected Returns.
  \textit{The Review of Financial Studies}, 29(1), 5--68.
  \url{https://doi.org/10.1093/rfs/hhv059}

\item Hou, K., Xue, C., y Zhang, L. (2015). Digesting Anomalies: An Investment Approach.
  \textit{The Review of Financial Studies}, 28(3), 650--705.
  \url{https://doi.org/10.1093/rfs/hhu068}

\item Hou, K., Xue, C., y Zhang, L. (2020). Replicating Anomalies.
  \textit{The Review of Financial Studies}, 33(5), 2019--2133.
  \url{https://doi.org/10.1093/rfs/hhy131}

\item Jegadeesh, N., y Titman, S. (1993). Returns to Buying Winners and Selling Losers:
  Implications for Stock Market Efficiency. \textit{The Journal of Finance}, 48(1), 65--91.
  \url{https://doi.org/10.1111/j.1540-6261.1993.tb04702.x}

\item Ke, G., Meng, Q., Finley, T., Wang, T., Chen, W., Ma, W., Ye, Q., y Liu, T.-Y. (2017).
  LightGBM: A Highly Efficient Gradient Boosting Decision Tree.
  \textit{Advances in Neural Information Processing Systems}, 30, 3146--3154.

\item Kelly, B. T., Pruitt, S., y Su, Y. (2019). Characteristics Are Covariances: A Unified Model of
  Risk and Return. \textit{Journal of Financial Economics}, 134(3), 501--524.
  \url{https://doi.org/10.1016/j.jfineco.2019.05.001}

\item López de Prado, M. (2018). \textit{Advances in Financial Machine Learning}.
  Hoboken, NJ: John Wiley \& Sons.

\item López de Prado, M., Bailey, D. H., Borwein, J. M., y Zhu, Q. J. (2014). Pseudo-Mathematics
  and Financial Charlatanism: The Effects of Backtest Overfitting on Out-of-Sample Performance.
  \textit{Notices of the American Mathematical Society}, 61(5), 458--471.
  \url{https://doi.org/10.1090/noti1105}

\item McLean, R. D., y Pontiff, J. (2016). Does Academic Research Destroy Stock Return
  Predictability? \textit{The Journal of Finance}, 71(1), 5--32.
  \url{https://doi.org/10.1111/jofi.12365}

\item Newey, W. K., y West, K. D. (1987). A Simple, Positive Semi-Definite,
  Heteroskedasticity and Autocorrelation Consistent Covariance Matrix.
  \textit{Econometrica}, 55(3), 703--708.
  \url{https://doi.org/10.2307/1913610}

\item Novy-Marx, R. (2013). The Other Side of Value: The Gross Profitability Premium.
  \textit{Journal of Financial Economics}, 108(1), 1--28.
  \url{https://doi.org/10.1016/j.jfineco.2013.01.003}

\item Novy-Marx, R., y Velikov, M. (2016). A Taxonomy of Anomalies and Their Trading Costs.
  \textit{The Review of Financial Studies}, 29(1), 104--147.
  \url{https://doi.org/10.1093/rfs/hhv063}

\item S\&P Dow Jones Indices (2026). \textit{SPIVA U.S. Scorecard: Year-End 2025}.
  Tabla 1a; datos a 31 de diciembre de 2025.
  \url{https://www.spglobal.com/spdji/en/documents/spiva/spiva-us-year-end-2025.pdf}

\item Sharpe, W. F. (1991). The Arithmetic of Active Management.
  \textit{Financial Analysts Journal}, 47(1), 7--9.
  \url{https://doi.org/10.2469/faj.v47.n1.7}

\item Shumway, T. (1997). The Delisting Bias in CRSP Data.
  \textit{The Journal of Finance}, 52(1), 327--340.
  \url{https://doi.org/10.1111/j.1540-6261.1997.tb03818.x}

\item White, H. (2000). A Reality Check for Data Snooping.
  \textit{Econometrica}, 68(5), 1097--1126.
  \url{https://doi.org/10.1111/1468-0262.00152}

\end{list}
